
\subsection{Continue ordinary events on the tax-adjusted book}

The additional tax payable must survive when ordinary business resumes. A
later return refund or lease payment has to use the cash left after tax
payments, and its journal must leave the tax balances intact unless that
method explicitly changes them. The continuation profile therefore retains
all 26 accounts, the return and lease memories, the assessment result and one
ordered event history. It extends the typed state and memory schemas directly;
it does not remove the tax payable to fit the older book.

The selected return methods instantiate the existing monetary definitions over
the complete account registry. Their memory operations still settle the return,
create the recovered lot or dispose of that lot, while preserving the additional
balances. The selected lease methods retain their original component equations
and map the resulting journal into the same registry. Tax expense and deferred
balances are unchanged by these ordinary methods. A separate, explicitly
instructed payment reduces income-tax payable and cash.

Assessment coverage constrains the next event. An ordinary event must occur on
or after the exclusive assessment endpoint and before the original financial
period ends. Otherwise it would change activity already used in the retained
tax calculation. A tax payment may occur on the accrual day, but cannot precede
an already completed operation. If assessment reaches the end of the declared
financial period, no ordinary continuation is possible without a separately
specified transition to the next opening.

\begin{samepage}
An earlier assessment provides a concrete same-period test. The author declares
the original two-year financial window, assesses the first interval through
31 December 2024 and places the later ordinary events on 1 January 2025.
The taxable expression of 2 and tax of $1/2$ are explicit hypothetical
assumptions. Paying that tax, settling the return, selling the recovered lot
and paying the lease use one cash balance:
\[
 500-\tfrac12-48+30-121=\tfrac{721}{2}.
\]
\end{samepage}
After the accompanying depreciation, recorded profit is $9/2$, assets are
$1401/2$, liabilities are 104 and equity is $1193/2$. Operating cash is
$-59/2$ and financing cash is $-110$. Direct cash, the indirect bridge,
the balance sheet and the equity movement reconcile exactly. The return's
later tax adjustment of $-24$ remains dated and identified for a subsequent
assessment; this continuation does not silently fold it into the earlier tax.

Payment order becomes consequential when cash is scarce. With opening cash
$485/4$, the lease payment of 121 alone leaves $1/4$. Paying tax of $1/2$
first leaves $483/4$, which cannot fund that lease payment. Reversing the order
allows the lease payment but leaves too little for tax settlement. Each request
stops at its first failed operation and retains the complete successful prefix.
It neither restores spent cash nor supplies financing automatically.

The result preserves the complete original report, loss state and conditional
tax calculation alongside the subsequent journals and gross cash movements.
Applying it recomputes the ordered prefix and compares the whole result.
Unknown findings remain unresolved, and repeated identities, changed balances
and backdated activity are refused. These checks establish selected same-period
execution under the declared assumptions. Further tax assessment, financial-year
carry and actual accounting admission remain separate work.
